\chapter{Tableaux de Butcher}\label[appendix]{app-butcher}

Une méthode de Runge--Kutta à $s$ étages s'écrit $y_{n+1}=y_n+h\sum_{i=1}^{s}b_ik_i$ avec $k_i=f\bigl(t_n+c_ih,\;y_n+h\sum_{j<i}a_{ij}k_j\bigr)$. Ses coefficients $c_i$, $a_{ij}$ et $b_i$ se présentent dans un \emph{tableau de Butcher} : la première colonne contient les $c_i$, le bloc central la matrice $A=(a_{ij})$ et la dernière ligne les poids $b_i$.

\bigskip
\begin{center}
\renewcommand{\arraystretch}{1.25}
\begin{tabular}{c@{\hspace{3cm}}c}
$\begin{array}{c|cc}
  0 &     &     \\
  1 & 1   &     \\ \hline
    & \tfrac12 & \tfrac12
\end{array}$
&
$\begin{array}{c|cccc}
  0        &          &          &   &   \\
  \tfrac12 & \tfrac12 &          &   &   \\
  \tfrac12 & 0        & \tfrac12 &   &   \\
  1        & 0        & 0        & 1 &   \\ \hline
           & \tfrac16 & \tfrac13 & \tfrac13 & \tfrac16
\end{array}$\\[4pt]
\small Méthode de Heun & \small Méthode RK4 classique
\end{tabular}
\end{center}

\chapter{Exemple de code : la méthode RK4}\label[appendix]{app-code}

La fonction suivante intègre un système $y'=f(t,y)$ avec la méthode de Runge--Kutta d'ordre~4 ; elle correspond à la définition du \cref{sec-classiques}.

\begin{lstlisting}[caption={Méthode RK4 en Python (pas constant).},label={lst-rk4}]
def rk4(f, t0, y0, T, N):
    """Integre y' = f(t, y) sur [t0, T] avec N pas de taille h."""
    h = (T - t0) / N
    t, y = t0, y0
    ys = [y0]
    for _ in range(N):
        k1 = f(t, y)
        k2 = f(t + h / 2, y + h / 2 * k1)
        k3 = f(t + h / 2, y + h / 2 * k2)
        k4 = f(t + h, y + h * k3)
        y = y + h / 6 * (k1 + 2 * k2 + 2 * k3 + k4)
        t = t + h
        ys.append(y)
    return ys
\end{lstlisting}
